\documentclass[titlepage]{article}

\usepackage{
	geometry,
	multicol,
	pdfpages
}

\geometry{
	letterpaper,
	margin=1in
}

\setlength{\columnsep}{0.25in}

\title{
	\textbf{
		CSCE 221 --- 200 \\
		Programming Assignment 4 \\
		Final Report
	}
}

\author{
	James Corder Guy \\
	Nathan Powell
}

\date{
	\today
}

\begin{document}
	\includepdf[pages={-}]{./coverpage.pdf}
	\maketitle
	\begin{multicols}{2}
		\section{Introduction}
                This report analyzes the implemention of a fully-functional graph data structure, including a searching algorithm and a minimum spanning tree algorithm. [add more once analysis is finished.]
		\section{Implementation Details}
			\subsection{Graph}
                        The graph was implemented using a map of vertices, a master map of edges, and an adjacency edge map for each vertex. The vertex map keys the vertexs by their identifiers, and each element contains a pointer to its respective vertex. Each vertex is assigned a unique identifier number corresponding to the number at which it was inputted, and with a list of edges that connect to it. Both edge maps key the edges by the vertex identifier pairs that they connect, and each element contains a pointer to its respective edge. Each edge is identified by the two vertices it connects, and its weight. All edges are directed; in order to insert an undirected edge, two edges going in opposite directions to the same two vertices are inserted.
			\subsection{Input/Output}
                        The graph pulls its initial vertices and edges from a file, reading in first the number of vertices, then the number of edges. It then reads in the next n lines, where n is the number of vertices, and initializes them as vertices in the vertex container. The graph then reads in the remaining lines as edges, with each line in the order of source vertex, destination vertex, and weight.
			\subsection{Breadth-First Search}
			\subsection{Minimum Spanning Tree}
                        Kruskal's Algorithm was used to find the minimum spanning tree of the input graph. The edges in the minimum spanning tree were copied into a parent map. To start, the algorithm read all the edges in the graph into a multimap, keying the edges by weight. Next, the first and last keys in the multimap are checked; if they are the same, then every edge in the map must have the same weight; therefore, every spanning tree is a minimum spanning tree. If not, the algorithm continues and sets up containers known as clusters; these represent the vertices linked together by Krusr ekal's algorithm that will form the minimum spanning tree. Another container that holds pointers to each cluster is initialized for constant access to each cluster. Every vertex is then inserted into its own cluster, and the algorithm begins to iterate through every edge in the multimap, comparing the vertices it connects; if they are not in the same cluster, that edge is added to the parent map, and the smaller cluster is merged into the larger cluster, though the now empty smaller cluster is not deleted; doing so would reindex the remaining clusters, invalidating all of the pointers to the clusters. This process repeats until there are n-1 edges in the parent map, with n being the number of vertices in the graph.
		\section{Theoretical Analysis}
			\subsection{Graph}
			\subsection{Input/Output}
			\subsection{Breadth-First Search}
			\subsection{Minimum Spanning Tree}
                        Kruskal's Algorithm should run in O((n+m)logn) time, because each vertex will be merged at most log(n) times, resulting in O(nlogn) time, because smaller clusters are always merged into larger clusters, and there will be at most m removals from the edge multimap (because there are m edges in total), resulting in O(mlogn) time, because each union operation on the edge's vertices takes O(logn) time. Therefore, the total running time of Kruskao's Algorithm is O((n+m)logn) time.

		\section{Experimental Analysis}
			\subsection{Testing Hardware}
                        Kruskal's Algorithm:
                          OS: Linux Mint 17.3 Cinnamon  64-bit
                          Kernel: 3.19.0-32-generic
                          Processor: Intel Core i5-2430M CPU @ 2.40 GHz x 2
                          Memory: 7.7 GB
                        Input Size:
                          2^(11) for complete, 2^(15) for mesh, 2^(14) for random
			\subsection{Results}
			\subsection{Big-O Constants}
		\section{Results and Discussion}
		\section{Team Contributions}
                Nathan was responsible for implementing the iterator operations, the insertion operations for vertex and edge, and the input operation, as well as Kruskal's Algorithm.
		\section{Conclusion}
	\end{multicols}
\end{document}
